%% exemplo-formulario.tex
%% Documento mínimo de teste — linha de formulários CRP-SP.
%% Compilar 2x com LuaLaTeX (via podman) por causa de \pageref{LastPage}:
%%   sh ../docker/lualatex-podman.sh exemplo-formulario.tex
%% Fase 1: sem \DocumentMetadata (tagging PDF/UA-2 é fase 2).

\documentclass{crpsp-formulario}

\hypersetup{
  colorlinks     = false,
  allbordercolors = white,
  pdftitle  = {Exemplo de formulário CRP-SP},
  pdfauthor = {Conselho Regional de Psicologia de São Paulo},
}
\title{Exemplo de formulário CRP-SP}
\author{Conselho Regional de Psicologia de São Paulo}

\begin{document}
\Form
\section{Exemplo de formulário}

\textbf{Teste do sistema de campos auto-dimensionáveis.}
\vspace{\baselineskip}

% 6.1 — campo de linha cheia: a caixa chega à margem direita não importa
% o tamanho do rótulo (compare as duas linhas abaixo).
\campolinha{nome}{Nome completo:}
\campolinha{orgao}{Órgão expedidor do documento de identidade:}

\vspace{\baselineskip}
% 6.2 — linha rotulada multi-campo, pesos iguais (divisão igual)
\begin{camposlinha}
  \campo{banco}{Banco:}
  \campo{agencia}{Agência:}
  \campo{conta}{Conta:}
\end{camposlinha}

% 6.2 — pesos diferentes: o primeiro campo fica com o dobro da largura
\begin{camposlinha}
  \campo[2]{endereco}{Endereço:}
  \campo[1]{numero}{Nº:}
\end{camposlinha}

\vspace{\baselineskip}
% 6.3 — campo de largura fixa dentro de prosa justificada
\justifying
Eu, \campofixo{decl-nome}{}{18em}, portador do RG
\campofixo{decl-rg}{}{9em}, declaro para os devidos fins.
\vspace{\baselineskip}

% 6.4 — datas e checkboxes
Data de nascimento: \campodata{nasc}
\vspace{0.5\baselineskip}

\begin{caixaslista}
  \caixaitem{opt-clinica}{Clínica}
  \caixaitem{opt-escolar}{Escolar/Educacional}
  \caixaitem{opt-social}{Social}
\end{caixaslista}
\vspace{\baselineskip}

% moldura + alias de compatibilidade (\preenctext legado, largura fixa)
\begin{framed}
  \preenctext{14}{compat}{20em}{Compatibilidade (\textbackslash preenctext): }
\end{framed}
\vfill

\campodataextenso{fim}
\vspace{2\baselineskip}

\blocoassinatura{Assinatura}

\end{document}
